\chapter{Antecedentes y estado del arte}

\label{chap:relacionados}

El análisis automático de ECG exige decidir qué información recibe el modelo y qué se espera que reconozca a partir de ella. Un registro completo permite estudiar el ritmo y comparar derivaciones, mientras que un latido aislado concentra la atención en su morfología. También cambia la referencia necesaria para evaluar la predicción, pues reconocer una alteración visible en el trazado y clasificar un paciente según su diagnóstico clínico son tareas con distinto alcance. Estas diferencias permiten ordenar los antecedentes y situar la pregunta de este trabajo, que estudia cómo aprovechar pocos ECG etiquetados para clasificar imágenes de un latido.

\section{Clasificación a partir de señales e imágenes}

\label{sec:antecedentes_imagen}

La señal digital constituye el punto de partida habitual porque conserva las muestras de amplitud y su evolución temporal. Antes de pasar a las imágenes, interesa examinar qué han permitido clasificar esas señales y cómo se han evaluado las predicciones. A partir de ahí se puede estudiar qué información permanece disponible al dibujar el trazado y qué dificultades aparecen cuando los ejemplos de la alteración que se quiere reconocer son escasos.

\subsection{Análisis de señales digitales}

La información que necesita un clasificador depende de la pregunta que se le plantea. Para estudiar el ritmo hacen falta varios latidos y sus intervalos, mientras que rasgos como la anchura del complejo QRS o la elevación del segmento ST pueden examinarse dentro de un latido. La duración del tramo y las derivaciones disponibles delimitan así qué puede observar el modelo, incluso antes de elegir su arquitectura. \cite{kligfield2007}

Ribeiro y colaboradores abordaron la detección de seis alteraciones del ritmo y de la conducción a partir de ECG de doce derivaciones del estudio CODE. Su red residual produjo predicciones por registro, entre ellas la del bloqueo de rama derecha, y se evaluó en ECG reservados para prueba con anotaciones de cardiólogos como referencia. Este tipo de tarea permite reconocer varias alteraciones en un mismo examen, pero su evaluación sigue vinculada a la definición de las etiquetas y a los registros elegidos para la prueba. \cite{ribeiro2020}

Para comparar arquitecturas bajo unas mismas condiciones, es necesario mantener estables esas etiquetas y los pacientes de evaluación. PTB-XL ofrece anotaciones de diagnóstico, forma y ritmo, junto con particiones por paciente que Strodthoff y colaboradores utilizaron para comparar redes como ResNet e Inception. Aunque este protocolo permite estudiar las diferencias entre modelos dentro de una muestra común, comprobar si se mantienen en otro centro exige ampliar la evaluación. El reto PhysioNet de 2020 incorporó datos ocultos de una institución ausente del entrenamiento y algunos métodos destacados en validación obtuvieron puntuaciones inferiores en la prueba final. La procedencia de los ECG importa, por tanto, junto con la separación de los pacientes, porque entre centros pueden variar la población, la adquisición y la obtención de las etiquetas. \cite{wagner2020,strodthoff2021ptbxl,perezalday2020challenge}

Hasta aquí, la predicción se refiere al registro completo. Al centrar la tarea en un latido, la referencia debe localizarlo e identificarlo dentro de la señal. He y colaboradores utilizaron anotaciones de este tipo en MIT-BIH, separaron los pacientes de entrenamiento y prueba y examinaron el método también en INCARTDB. Localizar los límites de las ondas exige anotaciones más detalladas, como las de la QT Database que Jimenez-Perez y colaboradores emplearon para evaluar la detección y la delimitación temporal de los componentes P, QRS y T. Estos trabajos muestran cómo la pregunta de clasificación determina el nivel de detalle exigido a la referencia, desde una etiqueta por registro hasta los límites de una onda. \cite{mitbih,incartdb,he2018pyramid,qtdb,jimenezperez2021delineation}

Disponer de muchas señales tampoco implica disponer de la referencia que necesita cada aplicación. MIMIC-IV-ECG vincula los registros con información clínica, y McKeen y colaboradores utilizaron señales de esta base y del reto PhysioNet de 2021 para preentrenar ECG-FM de forma autosupervisada. Al adaptar el modelo, estudiaron tareas con etiquetas de interpretación del ECG y con mediciones clínicas de la fracción de eyección ventricular izquierda, cuya evaluación quedó limitada al subconjunto que contaba con la medición correspondiente. El preentrenamiento permite aprovechar señales sin esa referencia específica, aunque la disponibilidad de casos etiquetados vuelve a delimitar la adaptación y la evaluación de la tarea final. \cite{gow2023mimicecg,reyna2022challenge,mckeen2025ecgfm}

En Brugada, además, importa la relación temporal entre el ECG y la referencia. Călburean y colaboradores estudiaron tanto el reconocimiento del patrón de tipo 1 en registros obtenidos durante una prueba con ajmalina como la predicción de su aparición a partir del ECG basal. Reconocer un patrón presente y anticipar un desenlace establecido en otro registro exige interpretar de forma distinta la predicción, aunque ambas tareas se relacionen con la misma alteración. La comparación entre derivaciones individuales y su combinación también mostró que la información puesta a disposición del modelo condiciona el resultado. \cite{calburean2024brugada}

Esta distinción es relevante para el presente TFG porque el patrón de Brugada puede ser intermitente y la valoración clínica puede apoyarse en antecedentes, otras derivaciones o pruebas de provocación. Un único latido recoge solo parte de esa información. Por ello, las etiquetas morfológicas del simulador permiten evaluar rasgos del trazado, mientras que la etiqueta clínica del HUCA permite estudiar la concordancia de la predicción con la clasificación del paciente. Una vez delimitado qué se evalúa en cada caso, queda por examinar cómo representar ese latido para que el modelo pueda interpretarlo. \cite{zeppenfeld2022,wilde2023scbt}

\subsection{Análisis de imágenes}

Al convertir la señal en una imagen, el tiempo y la amplitud pasan a expresarse mediante la posición del trazo. La morfología seguirá siendo visible en la medida en que el recorte, la escala y la resolución permitan distinguir los rasgos de interés. Cuando se dibuja una señal numérica es posible controlar estos factores, mientras que en una fotografía o un escaneo también intervienen la presentación del documento y las condiciones de captura. \cite{kligfield2007,reyna2024challenge}

El reto PhysioNet de 2024 evaluó por separado la reconstrucción de la señal y la clasificación directa de los píxeles como dos formas de aprovechar el documento. Para estudiar la segunda vía, Gliner y colaboradores compararon señales y trazados renderizados en la detección de fibrilación auricular e introdujeron variaciones de perspectiva y fondo que simulaban una captura con el móvil. Su experimento mostró que la representación visual podía conservar información útil para reconocer esa alteración. \cite{reyna2024challenge,gliner2020}

La comparación de Sangha y colaboradores amplió la pregunta a varias alteraciones del ritmo y de la conducción. A partir de las señales de CODE, ya utilizadas por Ribeiro, generaron imágenes con distintas disposiciones de las doce derivaciones y evaluaron los modelos también en PTB-XL. El rendimiento comparable entre señales e imágenes en las alteraciones estudiadas respalda el interés de la representación visual, aunque su utilidad para otra tarea depende de que el rasgo que se busca continúe siendo distinguible. En Brugada, esto dirige la atención a las derivaciones precordiales derechas y a la morfología del punto J, el segmento ST y la onda T. \cite{ribeiro2020,sangha2022,zeppenfeld2022,wilde2023scbt}

El TFG emplea un latido central de V1 con un recorte y una escala controlados, de modo que Gemma~4 y ResNet18 reciben el mismo trazado de prueba. La comparación puede centrarse así en cómo aprovechan los ejemplos etiquetados disponibles. Sin embargo, elegir una representación común no resuelve la dificultad de reunir suficientes ejemplos de una enfermedad poco frecuente para aprender la tarea, que es el problema que motiva este estudio. \cite{saleh2026brugada}

\subsection{Aprendizaje con pocos casos de Brugada}

Una forma de trabajar con pocos casos consiste en partir de lo aprendido con otros datos y utilizar los ECG etiquetados para ajustar el modelo a Brugada. Saleh y colaboradores compararon estrategias de preentrenamiento supervisado y autosupervisado, junto con técnicas de sobremuestreo, en la clasificación de ECG asociados al síndrome. Este antecedente sitúa la escasez de ejemplos como una condición que debe abordarse al diseñar el aprendizaje. El presente trabajo estudia la posibilidad de utilizar esos casos como ejemplos de referencia durante la consulta de un modelo de visión y lenguaje, manteniendo fijos sus pesos. \cite{saleh2026brugada}

\clearpage

\section{Aprendizaje en contexto con modelos de visión y lenguaje}

\label{sec:antecedentes_icl}

El aprendizaje en contexto permite orientar una nueva tarea mediante instrucciones y ejemplos incluidos en la propia consulta. En el ámbito multimodal, Flamingo mostró cómo un modelo preentrenado podía interpretar secuencias de imágenes y texto y responder a consultas a partir de ejemplos intercalados. Los casos disponibles pasan así a formar parte de la entrada, por lo que su selección resulta relevante para estudiar el comportamiento del modelo. El funcionamiento de esta técnica se desarrolla en el \hyperref[sec:fundamentos_icl]{apartado de fundamentos}. \cite{alayrac2022}

La aplicación de modelos de visión y lenguaje a los ECG también se ha abordado mediante entrenamiento especializado. Liu y colaboradores presentaron ECGInstruct, un conjunto de pares de imágenes e instrucciones empleado para entrenar PULSE, y evaluaron sus capacidades mediante ECGBench. Mientras que este enfoque incorpora conocimiento electrocardiográfico durante el entrenamiento, el ICL permite estudiar cómo responde un modelo generalista cuando recibe ejemplos de la tarea durante la inferencia. \cite{liu2026}

La posibilidad de utilizar demostraciones visuales para analizar ECG ya cuenta con antecedentes cercanos a este TFG. Camba y colaboradores las emplearon para identificar patrones asociados al bloqueo de rama izquierda y a Brugada, incorporando conceptos clínicos intermedios antes de la clasificación. Costa y colaboradores estudiaron cómo influían la representación visual y las instrucciones, y compararon modelos de visión y lenguaje con métodos especializados en un escenario de escasez de datos de Brugada. Estos trabajos sitúan la aportación del TFG en la comparación de Gemma~4, ejecutado localmente con sus pesos abiertos, y ResNet18 al clasificar rasgos morfológicos simulados, evaluar la transferencia entre datos sintéticos y clínicos y utilizar ejemplos del propio hospital. La muestra clínica utilizada es distinta de la de Costa y colaboradores, de modo que sus métricas y las del presente trabajo carecen de una base común para una comparación directa. \cite{camba2026icl,costa2026fewshot}

\section{Simulación y cambio de dominio}

\label{sec:antecedentes_adaptacion}

\label{sec:antecedentes_adaptacion_cnn}

Para examinar el reconocimiento de rasgos morfológicos conviene disponer de trazados cuyas variaciones se puedan controlar. Sayadi y colaboradores utilizaron componentes gaussianos en un modelo dinámico de ECG orientado al filtrado bayesiano. Esta representación permite variar la amplitud, la posición y la anchura de los componentes asociados a las ondas, y sirve como punto de partida para el simulador QRS/ST del TFG, desarrollado mediante la suma de componentes gaussianos. \cite{sayadi2010}

El control del trazado facilita el estudio inicial, aunque el paso a registros clínicos introduce diferencias visuales y de significado de las etiquetas. La adaptación de dominio aborda las diferencias entre las representaciones de los datos, sin garantizar que desaparezcan las discrepancias entre los criterios de clasificación. Gliner y colaboradores estudiaron esta adaptación entre imágenes digitales y fotografías de ECG mediante datos del dominio objetivo sin etiquetar, un antecedente de interés para examinar el uso de imágenes clínicas durante la adaptación. \cite{gliner2023}

Las estrategias consideradas en este TFG actúan de distintas formas sobre las representaciones. SimCLR permite aprenderlas mediante el contraste entre versiones aumentadas de una imagen, mientras que CORAL y MMD buscan aproximar las distribuciones de características mediante sus covarianzas o comparaciones basadas en kernels. DANN, por su parte, incorpora un clasificador adversarial de dominio para aprender representaciones que reduzcan la información sobre la procedencia de las imágenes. La evaluación de estas estrategias permite estudiar si el aprovechamiento de datos clínicos mejora la transferencia desde el simulador al HUCA, antes de comparar ese escenario con el uso de ejemplos etiquetados del propio hospital. \cite{chen2020simclr,sun2016deepcoral,long2015dan,ganin2016}
